% Historical source (removed by author request): references/_OLD_Übungsaufgaben/Übungsblatt 4 - Funktionen/Blatt4.tex:71
\begin{ExerciseBlock}
\ExerciseHeading{4}{Symmetrisierung von Funktionen}{functions_symmetrization}
Es sei $f$ eine beliebige Funktion auf einem zum Ursprung symmetrischen Definitionsbereich.
\begin{enumerate}[label=\alph*)]
	\item Rechnen Sie nach, dass 
\[\displaystyle g_g(x) = \frac{f(x)+f(-x)}{2}\] eine gerade Symmetrie hat.

	\item Rechnen Sie nach, dass 
\[\displaystyle g_u(x) = \frac{f(x)-f(-x)}{2}\] eine ungerade Symmetrie hat.

	\item Rechnen Sie nach, dass $\displaystyle f(x) = g_g(x) + g_u(x)$. Das bedeutet, dass sich jede Funktion $f(x)$, symmetrisch oder nicht, aus einem geraden und einem ungeraden Anteil zusammensetzt.

	\item Wählen Sie eine beliebige Funktion $f(x)$ aus und plotten Sie die Funktionen $f(x)$, $g_g(x)$ und $g_u(x)$.

\end{enumerate}

\end{ExerciseBlock}
\UseExerciseSolution{functions_symmetrization}
